\chapter{Early-init frame parameters and startup suppression}
\label{ch:early-init-frame}
\minitoc\label{early-init-frame-parameters-and-startup-suppression}
\index{default-frame-alist}\index{frame-inhibit-implied-resize}\index{early-init.el!frame parameters}\index{site-run-file}\index{vc-handled-backends}

% ── Opening signpost ──────────────────────────────────────────────
Chapter~\ref{ch:startup-lifecycle} introduced the two-file split and
noted that \texttt{early-init.el} runs before the first frame exists.
Several categories of startup overhead can only be avoided if the
relevant settings are applied before the first frame is created.
This chapter covers the four \texttt{early-init.el} mechanisms that
prevent frame resize flashes, suppress site-level interference, and
eliminate unnecessary version-control probes on every file open.

Four mechanisms in \texttt{early-init.el} prevent specific categories of startup overhead that cannot be avoided if set in \texttt{init.el}.

\section{\texttt{default-frame-alist}: prevent frame resize on creation}\label{default-frame-alist-prevent-frame-resize-on-creation}

Parameters pushed to \texttt{default-frame-alist} before the first frame is created take effect on that first frame without a resize cycle.

\begin{minted}{elisp}
;; early-init.el
(push '(menu-bar-lines . 0) default-frame-alist)
(push '(tool-bar-lines . 0) default-frame-alist)
(push '(vertical-scroll-bars) default-frame-alist)
(setq default-frame-alist '((width . 100) (height . 55)))
\end{minted}

If these were set after the frame exists (in \texttt{init.el}), Emacs would create the frame with default parameters and then resize it, causing a visible flash and extra computation.

\section{\texttt{frame-inhibit-implied-resize t}}\label{frame-inhibit-implied-resize-t}

Setting a font whose metrics differ from the default causes Emacs to resize the frame to maintain the same number of rows and columns (the "implied resize" behaviour). During early-init, the font is set before the frame exists, so this resize is normally free, but on some platforms it still triggers a synchronous geometry negotiation with the window manager. Setting this inhibits that:

\begin{minted}{elisp}
;; early-init.el
(setq frame-inhibit-implied-resize t)
\end{minted}

\section{\texttt{site-run-file nil}}\label{site-run-file-nil}

Emacs loads \texttt{/etc/emacs/site-start.el} (or the OS equivalent) before \texttt{init.el}. On shared systems this file may load packages, set variables, or add to \texttt{load-path} in ways that conflict with \texttt{straight.el} or shadow packages. Silencing it:

\begin{minted}{elisp}
;; early-init.el
(setq site-run-file nil)
\end{minted}

\section{\texttt{vc-handled-backends '(Git)}}\label{vc-handled-backends-git}

On every \texttt{find-file}, Emacs probes the file against every registered VC backend (Git, SVN, Mercurial, Bazaar, etc.) to determine file state. Restricting to Git eliminates all non-Git probes:

\begin{minted}{elisp}
;; early-init.el
(setq vc-handled-backends '(Git))
\end{minted}

Measured saving: \(\sim\)20--30 ms on startup and on every \texttt{find-file} call.

% ── Mental model summary ──────────────────────────────────────────
\section{Mental model}
The central theme of these settings is that they must take effect before the first frame is created or the package system initialises; placing them in \texttt{init.el} is too late.  Setting \texttt{default-frame-alist} entries (font, dimensions, toolbar/scrollbar visibility) before frame creation avoids a visible resize-and-flicker cycle.  \texttt{frame-inhibit-implied-resize} goes further by preventing the synchronous geometry negotiation that Emacs would otherwise perform whenever fonts or fringes change.

Two additional early settings improve startup hygiene.  \texttt{site-run-file} set to \texttt{nil} silences \texttt{/etc/emacs/site-start.el}, whose system-level configuration may conflict with \texttt{straight.el} or other user choices.  \texttt{vc-handled-backends} restricted to \texttt{'(Git)} eliminates probes for SVN, Mercurial, and Bazaar on every \texttt{find-file}, saving 20--30\,ms per file visit.  Together, these four \texttt{early-init.el} settings remove visible flicker, avoid redundant work, and prevent external configuration from interfering with the rest of the init process.
